% GOVERNANCE NOTE::LSPV3-WRITE-REPAIR::2026-10-12T08:08:28Z
\documentclass[11pt]{article}
\usepackage[margin=1in]{geometry}
\usepackage{xurl,hyperref}
\begin{document}
\section*{Governance note: LSPV-3 canonical-write repair}
\textbf{Change ID:} LSPV3-WRITE-REPAIR-20261012-0808.

\textbf{Evidence.} At GitHub preflight, the newest canonical LSPV-3 report
was \path{research/agents/lspv-3/2026-10-12_0532_LSPV-3_MINFRONT1.tex}.
It was substantive and had the required header and final sentinel.  The
enabled LSPV-3 automation showed later completed invocations but no newer
canonical report existed, covering at least two scheduled cycles.

\textbf{Classification.} \textbf{PROMPT / PROCEDURE CONTRADICTION.}  The
prompt requested a timestamped SHA-safe GitHub/Drive fallback but did not
state the canonical LSPV-3 path, mandatory GitHub readback/sentinel, or
same-run root-main ownership reconciliation required by \path{AGENTS.md}.

\textbf{Repair.} The automation schedule, enabled state, title, and
scientific target were preserved.  Its prompt received one appended
repository-write paragraph requiring the canonical LSPV-3 path, header and
sentinel, GitHub readback, current-SHA update of only the LSPV-3 region, and
no local/Drive success claim.  If an exact missing report body is available,
it may be retried once; missing mathematics may not be reconstructed.

\textbf{Rollback.} Remove only the appended paragraph beginning
``Repository write repair'' from automation
\texttt{6ac44a4d82648191bc8266c5523d92df}.  No schedule or scientific text
needs rollback.

\textbf{Status.} Repair applied and automation left enabled.  This note does
not assert that any missing scientific report existed or was recovered.

% END_OF_GOVERNANCE_NOTE::LSPV3-WRITE-REPAIR-20261012-0808
\end{document}
